\section{Programs, Executions and Traces}
\label{sec:model}

We model the point-to-point core of MPI with probes, completion calls and collectives.
\Cref{sec:implementation} lists what the implementation covers beyond it and what it
reports as unmodelled.

\subsection{Operations and programs}

A program runs on ranks $\Ranks=\{0,\dots,p-1\}$. Each rank issues a sequence of
\emph{operations}:
\begin{itemize}[leftmargin=*]
\item a \emph{send} $s$ with destination $\mathit{dst}(s)$, tag, communicator, size
  $\mathit{bytes}(s)$ and mode $\mathit{mode}(s)\in\{\mathsf{std},\allowbreak\mathsf{sync},\allowbreak
  \mathsf{buf},\allowbreak\mathsf{rdy}\}$, blocking or nonblocking;
\item a \emph{receive} $r$ with source $\mathit{src}(r)\in\Ranks\cup\{\ANY\}$, tag in
  $\mathbb{N}\cup\{\ANY\}$, communicator and capacity, blocking or nonblocking;
\item a \emph{probe} with the same envelope fields as a receive;
\item a \emph{completion call} (wait or test, single or over a set of requests);
\item a \emph{collective} with a name, a communicator and possibly a root.
\end{itemize}
Every operation $e$ has an \emph{issue} event $\issue(e)$ and a \emph{completion} event
$\compl(e)$. A blocking operation completes in its own call; a nonblocking one completes
in the completion call that returns its request. A matched pair $(r,s)$ has a \emph{match}
event $\mnode(r,s)$. Each call returns a \emph{result}: a request handle, a status (actual
source, tag, count and the received data), a flag, or an index.

\begin{definition}[Program]
A program $\prog=(\prog_i)_{i\in\Ranks}$ gives each rank a function $\prog_i$ from the
sequence of results it has received so far to its next operation, or to termination.
\end{definition}

A rank is thus deterministic given the results of its calls, and all nondeterminism comes
from the MPI library: which message a receive takes, when a request completes, what a test or
probe observes, and which request \mpicall{Waitany} returns. As the local-determinism axiom
does for shared memory~\citep{serbanuta12mcm}, the definition excludes control flow that
depends on clocks (including \mpicall{Wtime}, which the trace does not record), random
numbers, input or signals.

\subsection{Executions}

Let $\Bt\in\mathbb{N}\cup\{\infty\}$ be a buffering threshold. A send \emph{blocks at
$\Bt$} if its completion waits for its match:
\[
\blocks_\Bt(s) \iff \mathit{mode}(s)=\mathsf{sync} \;\lor\;
  \bigl(\mathit{mode}(s)\in\{\mathsf{std},\mathsf{rdy}\} \land \mathit{bytes}(s)>\Bt\bigr).
\]

\begin{definition}[Execution]
An \emph{execution} of $\prog$ at $\Bt$ is a finite sequence of issue, match and
completion events such that:
\begin{enumerate}[leftmargin=*,label=(E\arabic*)]
\item\label{E:po} each rank $i$ issues the operation $\prog_i$ yields on the results of
  its earlier calls, after the previous operation's call has returned;
\item\label{E:match} a match $\mnode(r,s)$ requires $r\accepts s$ and follows
  $\issue(r)$ and $\issue(s)$; no send and no receive is matched twice;
\item\label{E:fifo} (non-overtaking, sender side) if $s'$ precedes $s$ on the same rank,
  both to the same destination and communicator, and $r\accepts s'$, then $s'$ is matched
  before $\mnode(r,s)$;
\item\label{E:order} (non-overtaking, receiver side) if $r'$ precedes $r$ on the same rank
  and $r'\accepts s$, then $r'$ is matched before $\mnode(r,s)$;
\item\label{E:compl} $\compl(r)$ follows $\mnode(r,\cdot)$; $\compl(s)$ follows
  $\mnode(\cdot,s)$ if $\blocks_\Bt(s)$ and follows $\issue(s)$ otherwise; a ready send
  is erroneous unless a receive accepting it is posted and unmatched at $\issue(s)$;
\item\label{E:coll} a collective completes on a member only after the members its shape
  requires have issued it: all members for barrier-like collectives, the root for
  broadcast and scatter, all members at the root for reduce and gather, the lower ranks
  for scans;
\item\label{E:probe} a probe completes reporting a send $s$ that is issued and unmatched, such
  that a receive with the probe's envelope, posted at that point, would take $s$ under
  \ref{E:fifo} and \ref{E:order}.
\end{enumerate}
\end{definition}

Message delays are arbitrary: any interleaving that satisfies the rules is an execution.
A \mpicall{Waitany} may return any of its completed requests. A send that does not block
completes at once, so buffering is unbounded below the threshold. The threshold models an
implementation's eager protocol: the standard lets any standard-mode send block, which is
the case $\Bt=0$. An execution
\emph{deadlocks} if it cannot be extended and some rank has issued an operation that has
not completed.

\subsection{Traces and consistency}

\begin{definition}[Trace]
A \emph{trace} $\trace$ records, for each rank, the sequence of its calls with their
arguments and results. Its \emph{observed matching} $\Mobs$ pairs each completed receive
with the send it took.
\end{definition}

The tracer does not see match events. $\Mobs$ is reconstructed from the statuses:
per receiving rank, communicator and actual source, receives take that source's messages
in send order. \ref{E:fifo} makes this reconstruction exact.

A \emph{prefix} $\Hpre$ of $\trace$ chooses, for each rank, a prefix of the calls and returns
of its records, so it may keep a call without its return.

\begin{definition}[Consistency]
A program $\prog$ is \emph{consistent with $\trace$ on $\Hpre$} if, for every rank $i$
and every record $k$ of rank $i$ whose call is in $\Hpre$, $\prog_i$ applied to the results of
records $0,\dots,k-1$ of rank $i$ in $\trace$ yields the operation of record $k$. We write
$\Cons(\trace,\Hpre)$ for the set of such programs. The \emph{single-path program}
$\SP(\trace)$ issues the records of $\trace$ whatever its calls return.
\end{definition}

Consistency constrains a program only on the histories $\trace$ shows. It says nothing
about what a rank does after a call returns a result that $\trace$ does not record, for
example a different status. This is what makes predictions sound for the unknown program
behind the trace. The definition is the local-determinism axiom of the maximal causal
model~\citep{serbanuta12mcm}, stated for MPI results and relative to a prefix: every program
that produces $\trace$ is in $\Cons(\trace,\Hpre)$, so a claim about all of
$\Cons(\trace,\Hpre)$ is at least as strong as one about the programs that produce $\trace$.
The definition assumes that $\trace$ records every communication call of the run. The tracer reports calls outside the modelled fragment, and the analysis then
restricts itself to observed findings (\cref{sec:implementation}).

\subsection{Alternative matchings}

An \emph{alternative} is a partial injective map $\Malt$ from receives to sends with
$r\accepts\Malt(r)$. Its \emph{changed receives} are
$\Changed=\{r\in\mathrm{dom}\,\Malt \mid \Malt(r)\ne\Mobs(r)\}$. A \emph{steal} changes one
receive $r$ to a send $s$; if $s$ was taken in $\trace$ by a later receive $r_2$, $r_2$ loses its message.
A \emph{swap} exchanges the messages of two wildcard receives of one rank. By
\ref{E:fifo}, only the earliest message of each source that $r$ has not consumed can be
$\Malt(r)$, so a wildcard receive has at most one steal per source with a pending
message. The race sets of reachability testing~\citep{lei06reachability} and the racing
messages of the MAD debugger~\citep{kranzlmueller00thesis} have the same bound. A swap is
determined by its pair of receives (\cref{thm:cost}).
Two further kinds keep $\Malt=\Mobs$ and change one record instead. A \emph{probe
alternative} lets a probe report another send, and a \emph{completion-order alternative} lets
a \mpicall{Waitany} or \mpicall{Waitsome} return another request whose receive completes later in
$\trace$. The results below cover them with that record in the place of a changed receive
(\cref{app:conv}).
